\subsection{Row-independent actions of a two-row band}

\begin{theorem}[Proved: the row-swap obstruction opposite $R_3$]
\label{CP2:A-two-row-band-row-independent-right-zero}
Let $M=K^1$, where
$K=\{e_{ij}:i\in\{0,1\},\ 0\le j<c\}$ is a rectangular
band, $c\ge2$.  Let $N=R_3$ and let $S$ be a finite
nonempty left-$N$, right-$M$ biset whose right action
is independent of the row:
\[
 se_{0j}=se_{1j}\qquad(s\in S,\ 0\le j<c).
\]
If the rational second Betti number of the collage is
at most one, all its rational degree-two products vanish.
In particular it cannot have the rational cohomology
ring of $\mathbb{CP}^2$.  Neither boundedness of its
homology nor a restriction on $|S|$ is required.

For $c=3$, the source is catalog index $41$ of order seven.
Thus its row-independent connecting actions are excluded
from the eighteen-arrow budget $(7,3,8)$.
\end{theorem}
\begin{proof}
Write $R=\mathbb QM$, $e=e_{00}$, $f=e_{10}$ and
$q=1-e$.  Let $\sigma$ interchange the two rows of
$K$ and fix the identity.  The assumed equalities of
right actions imply $s\sigma(m)=sm$ for every $m\in M$.
Consequently $\sigma$ extends to an actual involution
of the collage by fixing $N$ and every connecting state.
We compute this actual involution on the bridge homology.

The explicit band resolution of
Lemma~\ref{CP2:A-band-spheres}, specialized to two rows,
is
\[
 0\longrightarrow P_2=(Re)^{c-1}
   \xrightarrow{d_2}P_1=Rq
   \xrightarrow{d_1}P_0=Re
   \xrightarrow{\varepsilon}\mathbb Q\longrightarrow0.
\]
Its maps are
\[
 d_1(wq)=w(f-e),\qquad
 d_{2,j}(we)=w(e_{0j}-e),\quad 1\le j<c.
                                                               \tag{1}
\]
Define $\sigma$-semilinear maps on these projectives by
\[
 F_0(u)=\sigma(u),\qquad
 F_1(wq)=-\sigma(w)q,\qquad
 F_{2,j}(u)=-\sigma(u).                                      \tag{2}
\]
These formulas are well-defined.  The ideal $Re=Rf$
has basis $e,f$ and is preserved by $\sigma$.
Also the kernel of $R\to Rq$, $w\mapsto wq$, is
$Re$; since that kernel is $\sigma$ invariant, $F_1$
does not depend on the chosen representative $w$.

The maps in (2) commute with (1).  For $d_1$, both
composites on $wq$ are $-\sigma(w)(f-e)$.
For $d_{2,j}$, both composites on $we$ are
\[
 -\sigma(w)(e_{1j}-f),
\]
using $f(e_{0j}-e)=e_{1j}-f$ and
$e_{1j}q=e_{1j}-f$.
They lift the identity on the augmented trivial module,
and square to the identity.  The projective-resolution
comparison theorem therefore identifies their action
with the actual row-swap action on the bar construction:
both are semilinear chain lifts of the same augmented
map, hence are chain homotopic after twisting the
source modules by $\sigma$.

Put $L=\mathbb Q[S]$, retaining its commuting left
$\mathbb QN$ action.  Tensoring (1) with $L$ gives
the natural left-$\mathbb QN$ complex
\[
 C_2=(Le)^{c-1}\longrightarrow C_1=Lq
                   \longrightarrow C_0=Le.
\]
The right-action hypothesis gives $Lf=Le$ with
$sf=se$ for every state, so the tensor differential
$C_1\to C_0$ is zero.  On these three terms the
involution induced by (2) is, respectively,
\[
 F_2=-1,\qquad F_1=-1,\qquad F_0=+1.                    \tag{3}
\]
For example $s\otimes e$ in degree zero is sent to
$s\otimes f$, which corresponds to $sf=se$;
the degree-two formula has the additional minus sign.
All comparisons and homotopies remain left-$N$ linear,
because they act on the projective $M$ resolution and
the two actions on $L$ commute.

Let $E=B(N,S,M)$ be the actual bridge.  Its rational
chains are computed by
\[
 \mathbb Q\otimes_{\mathbb QN}^{\mathbf L}C.
\]
The positive eigensubcomplex in (3) is exactly $Le$
in degree zero.  The trivial right $\mathbb QN$ module
has projective dimension at most one: if the two band
elements of $N$ are $u,v$, its resolution is
\[
 0\longrightarrow(1-u)\mathbb QN
  \xrightarrow{\ 1-u\mapsto v-u\ }u\mathbb QN
  \longrightarrow\mathbb Q\longrightarrow0.
\]
It follows that the positive eigenspace of $H_j(E;\mathbb Q)$
vanishes for $j\ge2$.  In particular $\sigma$ acts as
$-1$ on $H_3(E;\mathbb Q)$, irrespective of the
homology of the negative two-term summand.

By the band-sphere theorem, $BM$ is a wedge of
$c-1$ two-spheres and $BN$ is contractible.  The
two-block Mayer--Vietoris sequence, natural under
$\sigma$, therefore gives
\[
 H_4(BC;\mathbb Q)\cong H_3(E;\mathbb Q).
\]
Thus $\sigma$ acts as $-1$ on rational fourth homology
and cohomology of the collage.  Tensoring the same
resolution with the trivial right $R$ module shows
that $\sigma$ acts as $-1$ on $H^2(BM;\mathbb Q)$:
the middle term is zero and the degree-two term in
(3) is entirely negative.

It remains to use the second Betti-number hypothesis.
If $H^2(BC;\mathbb Q)=0$, the conclusion is immediate.
Otherwise choose a generator $\alpha$ of that one-dimensional
space.  If its restriction to $BM$ is nonzero, naturality
forces $\sigma^*\alpha=-\alpha$.  Therefore
\[
 \sigma^*(\alpha^2)=(\sigma^*\alpha)^2=\alpha^2,
 \qquad \sigma^*(\alpha^2)=-\alpha^2,
\]
so $\alpha^2=0$ over $\mathbb Q$.
If its restriction to $BM$ is zero, first collapse
the actual contractible corner $BN$.  This gives
the actual cofiber $Y=\operatorname{hocofib}(E\to BM)$
and a homotopy equivalence from the collage.
The class is then relative to $BM$ and lifts from
$Y/BM\simeq\Sigma E$, whose positive-degree products
are zero.  Its square again vanishes.  In a vector
space of dimension one this kills every product of
two second-degree classes, completing the proof.

For $c=3$ the literal band labeling is
$e_{ij}=1+3i+j$, with identity $0$.
The multiplication and all three semilinear chain
maps are independently checked in
\path{converse/scope18/two_row_band_source41_literal_binding.json}.
\end{proof}
